把一张正方形卡片转四分之一圈，放回桌上，看不出动过；把一个分子的原子换个编号顺序，它还是那个分子；把一段音频整体延后半秒，说的还是同一句话。这些场景共用一个模式：某些改动发生了，而我们关心的那部分没变。

这个模式平时被叫做``对称''，可这个词太像形容词了，听上去在描述图形好不好看。能拿来算的是那批改动本身——它们能连着做，能撤销，连做的次序有时还会影响结果。群这套语言把``哪些改动不改变本质''收集起来，当成一个能参与运算的对象。

本章的路线是四步：先从具体的对称动作提炼出群，再让群作用到外部对象上从而得到轨道，接着数清轨道，最后把轨道的要求翻译成对函数的约束。

\begin{center}
\begin{tikzpicture}[x=1cm,y=1cm,font=\footnotesize,>=stealth]
  \foreach \px/\ttl/\sub in {%
    0/{对称动作}/{转、翻、重排},
    3.3/{群}/{复合与可逆},
    6.6/{群作用}/{落到对象上},
    9.9/{轨道}/{哪些算同一个},
    13.2/{不变与等变}/{写进模型}}
  {
    \draw[black!55,fill=blue!8] (\px,0) rectangle (\px+2.6,1.1);
    \node[black!85] at (\px+1.3,0.72) {\ttl};
    \node[black!60,font=\scriptsize] at (\px+1.3,0.33) {\sub};
  }
  \foreach \px in {2.6,5.9,9.2,12.5}
    \draw[->,black!60] (\px+0.05,0.55) -- (\px+0.65,0.55);
\end{tikzpicture}

\emph{本章的四步。每一步都在把上一步的抽象重新接回可计算的东西：群把动作代数化，作用把群接回具体对象，轨道把``看起来一样''变成等价类，不变与等变把等价类变成对函数的约束。}
\end{center}

第一篇从正方形的八个对称动作出发，逼出封闭、结合、单位元、逆元四条要求，并说明每条挡住的是哪一种塌陷；顺带交代循环群、置换群与矩阵群三副常见面孔，以及非交换性为什么属于结构本身而不是技术细节。

第二篇让群落到对象上。子群是内部能独立运转的部分，陪集用子群把群切成等大的块并给出 Lagrange 定理，群作用则把抽象的群接回具体集合。轨道是``彼此能互相变过去''的等价类，稳定子是``按住不动''的那部分操作，轨道--稳定子定理说这两者的大小相乘恰好是群的大小。

第三篇把轨道数出来。轨道大小参差不齐，直接数很难；Burnside 引理靠一次数对子，把它换成``每个操作固定了多少对象''的平均，Pólya 的循环指标进一步让这些固定点数可以机械读出。这是去重计数的标准工具。

第四篇是本章与实践接得最紧的一篇。不变与等变是几何深度学习的两条设计原则：池化与读出层追求不变，卷积追求平移等变，图网络与点云模型追求置换等变，分子模型追求旋转平移等变。这些约束落到实处都是同一件事——参数共享，而共享方式由一个线性方程决定。同一篇也交代代价：对称假设一旦写进结构，就在所有输入上生效，选错时抹掉的信息拿不回来。

\begin{itemize}
\tightlist
\item \hyperref[sec:04-Discrete-Mathematics/06-Groups-and-Symmetry/01]{``对称操作引出群''}；
\item \hyperref[sec:04-Discrete-Mathematics/06-Groups-and-Symmetry/02]{``子群、陪集与群作用''}；
\item \hyperref[sec:04-Discrete-Mathematics/06-Groups-and-Symmetry/03]{``Burnside 与 Pólya：从轨道数到着色计数''}；
\item \hyperref[sec:04-Discrete-Mathematics/06-Groups-and-Symmetry/04]{``不变函数与等变映射：对称性怎样进入计算''}。
\end{itemize}

四篇按依赖排列，从具体问题进入也可以先打开最接近当前困惑的一篇，再沿篇内链接回补。若关心的是去重计数，第三篇需要第二篇的轨道语言；若从卷积网络或图网络进入，第四篇只需要第二篇末尾关于轨道与稳定子的那一段。

本章只搭起有限群与群作用的第一层语言。群的分类、Galois 理论不在范围内；把群系统地变成矩阵这件事留给\hyperref[ch:035]{``群表示论：对称的线性结构''}。读完能回答三个问题即可：允许哪些操作？它们作用在什么对象上？任务要保留轨道内的差异，还是主动把它抹掉？
